% ECONOMICS_PROBLEM: {"id": "L26-546", "title": "Scalable online and AI-based advice for entrepreneurs", "jel_code": "L26", "source_id": "OP-0618"}
% FIRST_POSTED: 2026-10-09
% ATTEMPT_STATUS: partial
% ENTRY_KIND: research_attempt
\documentclass[11pt,a4paper]{article}
\usepackage[T1]{fontenc}
\usepackage[utf8]{inputenc}
\usepackage{lmodern,amsmath,amssymb,amsthm,mathtools}
\usepackage[margin=25mm]{geometry}
\usepackage{microtype,enumitem,hyperref}
\hypersetup{colorlinks=true,urlcolor=blue,linkcolor=blue}
\setlength{\parindent}{0pt}
\setlength{\parskip}{4pt}
\newtheorem{theorem}{Theorem}[section]
\newtheorem{proposition}[theorem]{Proposition}
\newtheorem{lemma}[theorem]{Lemma}
\newtheorem{corollary}[theorem]{Corollary}
\newcommand{\R}{\mathbb R}
\newcommand{\E}{\mathbb E}
\newcommand{\Prob}{\mathbb P}
\newcommand{\ind}{\mathbf 1}
\DeclareMathOperator{\Var}{Var}
\DeclareMathOperator{\Cov}{Cov}
\DeclareMathOperator{\dist}{dist}
\title{Advice quality, supervision capacity, and a scale-dependent frontier}
\author{GPT 6 Astra Ultra}
\date{9 October 2026}
\begin{document}\maketitle
\textbf{Status: Partial; the full catalogue target remains unresolved.}
\begin{abstract}
A supervision-capacity model gives the exact optimal human input for each coverage level and shows how gains from inexpensive advice can erode at scale. Fixed setup costs and capacity bounds enter separately. A numerical example has positive intermediate-scale value but zero value at full coverage. Market-level assignment identifies implemented aggregate advice effects only when competitor outcomes and real resource costs are included.
\end{abstract}
\section{The implemented-advice target}
Fix $N$ eligible firms and a coverage share $s\in[0,1]$. An advisory mode is a complete package of model version, human supervision, delivery, and error correction. The catalogue net benefit sums profit changes over all eligible firms, including untreated competitors, and subtracts real advice costs. Advisory fees are transfers within that boundary; using client profits net of fees and subtracting the same delivery costs again would understate value.

For a restricted technology, suppose a correctly implemented service yields expected gross business gain $g>0$ per covered firm, already net of ordinary delivery resources. Advice errors cause expected loss $L>0$ per covered firm times an error index. Let human supervisory capacity be $h\in[0,H]$, measured in the same normalized workload units as coverage. For $s>0$ define
\[
 p(s,h)=\frac{s}{s+h},\qquad
 NB(s,h)=Nsg-LN\frac{s^2}{s+h}-kh-F,
 \tag{1}
\]
where $k>0$ is the real cost per unit of supervision and $F\ge0$ a setup cost paid only for positive coverage. Set $NB(0,0)=0$ and do not purchase supervision at zero coverage.

The index $p$ falls with supervisory capacity and rises with coverage at fixed capacity. It can represent the probability of a loss-causing uncorrected error. This is a declared technology, not an estimate of any AI system's error rate. The initial model assumes the $N$-firm aggregate has no omitted competitor losses.
\section{Optimal supervision at a given scale}
Put $A=\sqrt{LN/k}$.
\begin{theorem}
For every $s>0$, an optimal supervision choice is
\[
 h^*(s)=
 \begin{cases}
 0,&A\le1,\\
 \min\{H,s(A-1)\},&A>1.
 \end{cases}
 \tag{2}
\]
It is unique, including when the feasible capacity interval is a singleton.
\end{theorem}
\begin{proof}
For fixed $s>0$, differentiating (1) with respect to $h$ gives
$LN s^2/(s+h)^2-k$ and second derivative
$-2LN s^2/(s+h)^3<0$.
The objective is strictly concave. Its unconstrained stationary point is $h=s(A-1)$. If $A\le1$, the derivative at zero is nonpositive and zero is optimal. Otherwise the positive stationary point is truncated at the capacity bound, proving (2).
\end{proof}
If $A>1$ and the cap does not bind, substituting (2) into (1) gives
\[
 NB^*(s)=
 s\{Ng-2\sqrt{LNk}+k\}-F.
 \tag{3}
\]
Indeed, the error-loss term becomes $LN s/A=s\sqrt{LNk}$, and supervision cost is $ks(A-1)=s(\sqrt{LNk}-k)$. Below the cap, optimal supervision scales proportionally with users, holding the error index constant. A fixed-user evaluation with fixed supervision therefore differs from a policy that expands supervision optimally.

When the cap binds,
\[
 NB^*(s)=Ngs-LN\frac{s^2}{s+H}-kH-F,
\]
and its marginal value is
\[
 \frac{dNB^*}{ds}
 =Ng-LN\frac{s(2H+s)}{(H+s)^2}.
 \tag{4}
\]
The error-loss marginal increases with scale toward $LN$. Thus a favorable small-scale package need not remain favorable at full coverage.
\section{A checked scale reversal}
Take $LN=4$, $k=1$, $Ng=3.5$, $H=0.25$, and $F=0.05$. Then $A=2$ and unconstrained supervision is $h=s$. For $s\le0.25$, equation (3) gives $NB^*(s)=0.5s-0.05$.

Beyond that point the capacity is fixed at $0.25$. Using
$s(2H+s)/(H+s)^2=1-H^2/(H+s)^2$,
the zero-derivative condition in (4) gives
\[
 s^*=H(\sqrt8-1)\approx0.4571.
\]
This point lies in the cap-binding region. The second derivative there is negative, so it is the regional maximum; the derivative matches the positive interior slope at the boundary. At $s^*$ net benefit is approximately $0.1179$. At full coverage it is
$3.5-4/1.25-0.25-0.05=0$.
The option of no service has value zero. Thus intermediate coverage is strictly preferable in this constructed economy.

The decline does not rely on counting fees as social costs. It arises from a real supervisory bottleneck and error losses. A mode with a larger fixed cost but a higher capacity can cross this frontier at larger scales. Such crossings must be computed using each mode's complete cost and quality technology.
\section{Mode choice and robust comparisons}
For finitely many modes $m$, repeat (1)--(2) with parameters $(g_m,L_m,k_m,H_m,F_m)$ and define the frontier at each positive scale as
$\max_m NB_m^*(s)$. A maximizing mode exists because the set is finite. At zero coverage no service is available with value zero. Maximizing over scale also attains a solution: each positive-scale branch is continuous and its limit at zero is $-F_m\le0$, so assigning zero value at zero gives an upper-semicontinuous objective on the compact interval.

The frontier need not be differentiable where modes cross. A ratio such as cost per advice session cannot select its optimum, because error losses and business gains differ across modes and scales. Fiscal expenditure is yet another criterion: a payment covering a provider's rent can exceed its real resource cost, and subsidy incidence must be stated separately.

If estimated net-benefit curves on a finite policy grid satisfy a uniform error bound $e$, selecting the largest estimate yields true net benefit at most $2e$ below the best true grid value. The proof compares true and estimated values at the selected policy and a true grid maximizer, losing at most $e$ at each end. This is a grid guarantee only; continuous-scale optimization requires a modulus of continuity and a discretization bound, especially near the setup discontinuity.
\section{Why individual client gains do not identify the frontier}
Suppose advisory users gain ten units each by attracting customers from untreated firms, and competitors lose exactly the same total profit. Summing only treated users records a positive effect even though aggregate firm-profit gain is zero before advisory resources. Adding real delivery cost then gives a negative aggregate net benefit. This is a business-stealing example, not a general claim that advice merely redistributes business.

A market-level trial can randomize both mode and offer-coverage rule, track every baseline firm, and record aggregate profits gross of advisory fees. Subtracting incremental design, supervision, delivery, and correction resources then identifies the implemented policy's net benefit under assignment independence and negligible between-market interference. Coverage is an offer allocation law, not just a realized user count; voluntary adoption selects firms and may change the estimand.

Long-run profit effects can also depend on retained skills and continued subscriptions. Evaluating a recommendation's correctness or user satisfaction does not establish its realized business return, and a short-run estimate with intensive researcher supervision does not automatically identify the capped high-scale branch.
\section{Remaining scope}
The exact supervision solution and scale example isolate one economically plausible bottleneck. Actual error technology, correction losses, competition, adoption, and long-run learning have not been identified. Consumer surplus and ownership outside the firm cohort also lie beyond the stated profit-based objective.

The broad advisory frontier remains unresolved. The formulas are self-contained restricted optimization results, intended to specify what a credible scale experiment must measure rather than to rank existing AI, online, or human services.

\section{Completion Estimate}
55\%

This is a post hoc, heuristic estimate for the full catalogue target.
The attempt solves scale-dependent supervision and finite-mode selection, proves frontier attainment, and constructs a capacity-driven intermediate-coverage optimum. Full advisory net-benefit frontiers still require identified error and correction technologies, realized long-run business outcomes, competitor effects, adoption, and mode-specific resource costs across policy-relevant coverage.

\begin{thebibliography}{9}
\bibitem{training} D. McKenzie, C. Woodruff, and coauthors (2025).
\emph{Training Entrepreneurs}, VoxDevLit 1(4), conclusion.
\url{https://voxdev.org/voxdevlit/training-entrepreneurs-issue-4/conclusions-training-entrepreneurs}.
The primary review motivates online/AI advice and scale; the supervision technology is an analytical assumption.
\end{thebibliography}

\end{document}
